\silentchapter{Proposisjonar}{proposisjonar}

\begin{widebody}
\footnotesize
\textit{d} står for definisjon, \textit{a} for postulat, \textit{t} for teorem, \textit{o} for observasjon, \textit{i} for illustrasjon. Merket seier korleis ein proposisjon er vunnen, ikkje kva han er verd.
\end{widebody}

\vspace{6mm}

\mainprop{1}{}{Alt som finst, står. Alt som hender med det som står, er eitt og same forsøk.}
  \prop{1.1}{d}{Eit forsøk byter ein avgrensa del av det som står og lèt resten vere urørt.}
    \prop{1.11}{d}{Etter eit forsøk gjeld eitt av to ord: det som står, \textbf{held} — eller det spring attende, det \textbf{fell}.}
    \prop{1.12}{a}{Same forsøk, gjort om att på same ting, kan halde éin gong og falle ein annan.}
    \prop{1.13}{t}{Eit forsøk kan sjølv vere den avgrensa delen eit anna forsøk byter.}
  \prop{1.2}{d}{Eit forsøk som berre er tenkt, ikkje gjort, kallar vi ein hypotese.}
    \prop{1.21}{t}{Ein hypotese held, fell, eller står uprøvd. Eit fjerde utfall finst ikkje.}

\flowchapter{2 Vekt}{vekt}
\mainprop{2}{}{Ikkje alt som står, er like lett å nå med eit forsøk.}
  \prop{2.1}{d}{Vekta til noko som står, er talet på ulike kjeder av forsøk, frå ulikt utgangspunkt, som endar der.}
    \prop{2.11}{t}{Av 1.12 følgjer: høg vekt gjer ikkje neste forsøk sikrare. Ho gjev berre fleire vegar dit.}
    \prop{2.12}{o}{Låg vekt held stundom likevel. Vekta avgjer ikkje om noko står — berre kor lett det er å nå.}

\flowchapter{3 Endra grunn}{endra grunn}
\mainprop{3}{}{Kvart forsøk som held, endrar det som seinare forsøk møter.}
  \prop{3.1}{t}{Difor er grunnen eit forsøk møter, alltid summen — vekta, 2.1 — av dei forsøka som gjekk føre.}
    \prop{3.2}{t}{Same forsøk, møtt med ulik grunn, kan halde tidleg og falle seinare, utan at forsøket sjølv er endra.}

\flowchapter{4 Rekkevidd}{rekkevidd}
\mainprop{4}{}{Ikkje alle forsøk er opne for kven som helst.}
  \prop{4.1}{d}{Der forsøk på noko som står, er avgrensa til ein fast krins av utgangspunkt, kallar vi krinsen ein agent, og avgrensinga hans rekkevidd.}
    \prop{4.2}{t}{Ein agent kan difor aldri prøve alt som står — berre det rekkevidda hans når.}
    \prop{4.3}{t}{Sidan vekta til noko som står er summen av forsøk frå ulike utgangspunkt (2.1), og ein agent berre når nokre av dei, kan ingen agent åleine kjenne heile vekta til det han prøver.}

\flowchapter{5 Møte}{møte}
\mainprop{5}{}{Det same standande kan vere mål for fleire agentar sine forsøk på same tid.}
  \prop{5.1}{d}{Når to eller fleire agentar prøver det same standande på same tid, kallar vi det eit møte.}
    \prop{5.2}{t}{I eit møte kan forsøka krevje same avgrensa del.}
    \prop{5.3}{t}{Av 1.1 følgjer: då held berre eitt av dei. Dei andre fell — ikkje av eiga veikskap, men av at delen dei trong, alt var bytt.}

\flowchapter{6 Verknad mellom agentar}{verknad mellom agentar}
\mainprop{6}{}{Ein agent kan endre det ein annan agent møter, utan at dei nokon gong deler eit forsøk.}
  \prop{6.1}{t}{Av 3.1 og 4.1 følgjer: éin agent sin historie kan endre vekta ein annan agent finn, sjølv om ingen av rekkevidda deira når inn i den andre sin krins.}
    \prop{6.2}{o}{Ein agent kan difor verke på ein annan utan å vite det, og den andre kan bere verknaden utan å vite kven som gjorde han.}

\flowchapter{7 Forklaring og vurdering}{forklaring og vurdering}
\mainprop{7}{}{Å vite kvifor noko heldt, er ikkje det same som å vite kor mykje det er verdt at det heldt.}
  \prop{7.1}{d}{Ei forklaring er eit forsøk som peikar ut kva tidlegare forsøk avgjorde eit gjeve utfall.}
  \prop{7.2}{d}{Ei vurdering er ei rangering av det som står, etter vekt (2.1) — ikkje etter kva som avgjorde det.}
    \prop{7.3}{t}{Av 7.1 og 7.2 følgjer: to ting kan stå med lik vekt og ulik forklaring, eller lik forklaring og ulik vekt. Forklaring og vurdering møtest aldri av naudsyn.}

\flowchapter{8 Mynde}{mynde}
\mainprop{8}{}{Å avgjere kva forsøk som er opne for ein annan agent, er sjølv eit forsøk.}
  \prop{8.1}{t}{Ei rekkevidd (4.1) finst. Av mainprop 1 følgjer: alt som finst, står. Difor står ho òg — og kan sjølv vere den avgrensa delen eit forsøk byter.}
    \prop{8.2}{d}{Eit forsøk som byter ei rekkevidd, i staden for det rekkevidda gjeld, kallar vi eit myndeforsøk. Mynde er verknaden eit slikt forsøk har på ein annan agent.}
    \prop{8.3}{t}{Eit myndeforsøk krev difor korkje forklaring (7.1) eller vurdering (7.2) — som alle forsøk (1.11), er det avgjort ved om det held, ikkje ved om det er grunngjeve.}
    \prop{8.4}{o}{Den som prøver eit myndeforsøk på sin eigen rekkevidd, vender same rørsla mot seg sjølv.}

\flowchapter{9 Denne teksten}{denne-teksten}
\mainprop{9}{}{Denne traktaten er sjølv noko som står, bygd av forsøk — kvart utkast eit bytt av ein avgrensa del, resten urørt.}
  \prop{9.1}{t}{Vekta hans (2.1) er talet på lesarar sine forsøk på han — ikkje talet på setningar han inneheld.}
    \prop{9.2}{t}{Å lese er difor sjølv eit forsøk, av den slags 8.2 nemner. Det byter ikkje traktaten. Det byter lesaren sin rekkevidd.}
